% docs/texte/40_qa/tests.tex -- chapter "Tests and Validation" (Interferosim)
% Version: 0.6.0 (2026-11-25) -- notebook runs with nbval (N-4)
%   0.5.0 (2026-11-25) -- test certificate (tools/testzertifikat.py
%   1.1.0), tests/test_certificate.py (interferosim 1.0.5)
%   0.4.0 (2026-11-23) -- line coverage (pytest-cov), doctests,
%   property-based tests (Hypothesis), representative test for the
%   presentation; plain pytest also runs the doctests (interferosim 1.0.4)
%   0.3.1 (2026-11-23) -- table as longtable (overfull lines)
%   0.3.0 (2026-11-23) -- all eight test modules listed; selection by
%   name (-k), listing (--collect-only), uv run
%   0.2.0 (2026-10-12) -- running the tests
%   0.1.0: placeholder
% ================================================================
\label{ch:tests}
\index{tests}

Unit tests of the core package (pytest) and automated notebook runs
% 0.6.0: the notebook runs are set up (section \ref{sec:nbval})
%(nbval). The notebook runs follow in phase~5.
(nbval, section~\ref{sec:nbval}).

\section{Running the tests}

In the project folder with the project environment active
(\texttt{.venv\textbackslash Scripts\textbackslash Activate.ps1}):
\begin{lstlisting}[basicstyle=\ttfamily\footnotesize, frame=single]
pytest -q                                   # all tests, short output
pytest -v                                   # one line per test
pytest tests/test_sky.py -v                 # one module
pytest tests/test_sky.py::test_render_conserves_flux -v   # one test
pytest -k "hogbom or wiener" -v             # tests whose name matches
pytest --collect-only -q                    # list the tests, run none
pytest --cov                                # with line coverage (0.4.0)
pytest --junitxml=logs/tests/junit.xml      # test log for a certificate (0.5.0)
pytest --nbval-lax --current-env notebooks  # run all notebooks (0.6.0)
pytest tests/test_properties.py -v          # property-based tests only
.venv\Scripts\python.exe -m pytest -q       # without activating .venv
uv run pytest -q                            # via uv, without activating
\end{lstlisting}
% 0.3.0: the folder tests/ is set as testpaths in pyproject.toml, so a
% plain "pytest" finds all modules.
% 0.4.0: testpaths also contains src/interferosim (doctests)
%Without a path, pytest collects all modules in \texttt{tests/}; this
%folder is set as \texttt{testpaths} in \texttt{pyproject.toml}. With
Without a path, pytest collects all modules in \texttt{tests/} and,
because of \texttt{--doctest-modules}, the examples in the docstrings
of \texttt{src/interferosim/} (section~\ref{sec:doctests}); both
folders are set as \texttt{testpaths} in \texttt{pyproject.toml}. With
\texttt{-k}, a part of a test name or a boolean expression of names
selects the tests.

The tests are run before every commit; a release is only made when all
tests pass.

% ----------------------------------------------------------------
% 0.4.0: coverage, doctests, property-based tests, representative test
\section{Line coverage}\label{sec:coverage}
\index{coverage}

The line coverage of the package is measured with \texttt{pytest-cov}
(in the extra and dependency group \texttt{test}). The settings are in
\texttt{pyproject.toml}: \texttt{[tool.coverage.run]} measures the
package \texttt{interferosim}, \texttt{[tool.coverage.report]} lists the
missing lines and lets the run fail below 75\,\% (\texttt{fail\_under}).
\begin{lstlisting}[basicstyle=\ttfamily\footnotesize, frame=single]
pytest --cov                    # report in the terminal, fails < 75 %
pytest --cov --cov-report=html  # HTML report in htmlcov/index.html
\end{lstlisting}
Version~1.0.4 reaches 91\,\% (2596 statements, 237 not executed;
143 tests). Below average are only \texttt{plotting.imaging} and
\texttt{plotting.observation} (plot functions used by the notebooks) and
the command-line tool \texttt{sky.cli}.

\section{Doctests}\label{sec:doctests}
\index{doctest}

Selected functions and classes carry an \texttt{Examples} section in
their docstring, written as an interactive Python session. The examples
document the use and the expected result and are run as tests
(\texttt{--doctest-modules} in \texttt{addopts}, so with every
\texttt{pytest} call):
\begin{lstlisting}[basicstyle=\ttfamily\footnotesize, frame=single]
>>> from interferosim.constants import HI_21CM
>>> round(float(radio_velocity_kms(1420.0e6, HI_21CM)), 2)
85.64
\end{lstlisting}
Examples are in \texttt{constants.wavelength},
\texttt{sky.spectra} (radio velocity and its inverse),
\texttt{sky.components} (\texttt{PointSource}, \texttt{GaussianSource}),
\texttt{observation.geometry} (\texttt{gmst\_deg} at J2000,
\texttt{uvw}), \texttt{instrument.sensitivity} (radiometer equation,
visibility noise, coherence factor) and \texttt{signal.baseband} (image
rejection, quantisation efficiency, Van Vleck correction) -- 13 doctests.
Only the docstrings run:
\begin{lstlisting}[basicstyle=\ttfamily\footnotesize, frame=single]
pytest --doctest-modules src/interferosim
\end{lstlisting}
The values are rounded (\texttt{round}, \texttt{float}), so the output
does not depend on the way numpy prints its scalars.

\section{Property-based tests}\label{sec:hypothesis}
\index{Hypothesis}\index{property-based test}

\texttt{tests/test\_properties.py} uses Hypothesis. A test states a
property that must hold for \emph{all} inputs of a described range;
Hypothesis draws 100 random inputs per test and, if one fails, shrinks it
to the simplest counterexample. Checked properties:
\begin{itemize}
  \item frequency $\rightarrow$ radio velocity $\rightarrow$ frequency is
        the identity (all lines of \texttt{LINES});
  \item the visibility of every sky component (point, Gaussian, uniform
        disc) is Hermitian, $V(-u,-v) = V^*(u,v)$, as for every real
        brightness distribution; $|V| \le S$ with equality for a point
        source and $V(0,0) = S$;
  \item the amplitude of a Gaussian decreases along every ray in the
        uv-plane;
  \item \texttt{component\_from\_dict(c.to\_dict()) == c} (JSON round
        trip);
  \item the sky basis $(\vec s_0, \vec e_E, \vec e_N)$ is orthonormal
        and right-handed; $(u,v,w)$ is linear in the baseline and keeps
        its length;
  \item quantisation efficiency in $(0,1]$ and not decreasing with the
        number of bits; the Van Vleck correction is odd, monotonic,
        bounded and never smaller in magnitude than the 1-bit
        correlation; the image rejection ratio is not negative and even
        in the phase error;
  \item the coherence factor lies in $(0,1]$ and decreases with the
        phase rms.
\end{itemize}
A check with deliberately inserted errors shows that the two kinds of
tests complement each other: a wrong sign of $\vec e_E$ in
\texttt{sky\_basis} is found by two property tests at once, whereas a
visibility without its phase term is still Hermitian and bounded -- that
error is found by the doctest of \texttt{PointSource}, which expects the
imaginary part $-1.0944$.

\section{Representative test}\label{sec:representative-test}
\index{tests!representative}

For the final presentation:
\texttt{test\_visibility\_simulator\_agrees\_with\_iq\_simulation} in
\texttt{tests/test\_observation.py}.
\begin{lstlisting}[language=Python, basicstyle=\ttfamily\scriptsize, frame=single]
def test_visibility_simulator_agrees_with_iq_simulation():
    sky = SkyModel([PointSource("p", 0.0, 0.0, 2000.0)],
                   center=FieldCenter("p", 120.0, 45.0))
    site = Site.from_toml(CFG / "sites" / "wilhelmsburg.toml")
    dish = Antenna.from_toml(CFG / "antennas" / "dish_3m.toml")
    rec = Receiver.from_toml(CFG / "receivers" / "b210.toml")
    arr = Array(site, [Station.from_enu("A", site, 0, 0, 0, dish),
                       Station.from_enu("B", site, 200, 50, 0, dish)],
                receiver=rec)
    iq = IQSimulator(sky, arr, sample_rate_sps=2e4, seed=2)
    m = measure_visibilities(iq.simulate(20.0), iq, 20.0)
    vd_iq = VisData.from_measurements(m, iq.fc)
    assert abs(vd_iq.vis[0, 0] - 2000.0) < 5 * vd_iq.sigma[0, 0]
\end{lstlisting}

\paragraph*{Behaviour checked.} A point source of 2000\,Jy in the phase
centre is observed by two 3\,m dishes 206\,m apart with the receiver
profile of the B210. The closed simulation produces the noisy IQ samples
of both stations (geometric delay, receiver noise), the correlator
compensates the delay, stops the fringes and averages 20\,s. For a point
source in the phase centre the true visibility is known exactly: real,
$V = S = 2000\,\mathrm{Jy}$. The test checks that the measured
visibility agrees with it within five times the noise
$\sigma$ that the correlator itself estimates.

\paragraph*{Why this test.} It joins the two ways from the sky to the
visibilities that the design (chapter~\ref{ch:basic-design}) keeps
apart: the closed simulation as with real hardware (\texttt{signal}) and
the known analytic answer. It passes only if the whole chain works
together -- sky model, array geometry, delay model, IQ simulation,
correlator and \texttt{VisData} -- and it uses the real configuration
files. The comparison is against the noise estimate instead of a fixed
tolerance, so it also checks that the error bars are right. With a fixed
seed it is reproducible and runs in about a second.

\paragraph*{Bugs it would catch.}
\begin{itemize}
  \item a wrong sign or unit (seconds vs.\ samples) of the geometric
        delay or its compensation: the phase runs away, the real part
        drops far below 2000\,Jy;
  \item missing or wrong fringe stopping: the visibility rotates during
        the 20\,s and averages towards zero;
  \item a wrong flux scale (SEFD, Jy-to-K factor, normalisation of the
        correlator): the amplitude is off by a factor;
  \item a wrong noise estimate: a too small $\sigma$ makes the test
        fail, a too large one would show up as unrealistic error bars in
        the notebooks.
\end{itemize}

% ----------------------------------------------------------------
% 0.6.0: notebook runs
\section{Notebook runs (nbval)}\label{sec:nbval}
\index{nbval}\index{notebooks!tests}

The plugin nbval (in the extra and dependency group \texttt{test}) runs
every notebook from the first to the last cell in a fresh kernel and
reports each cell as a test:
\begin{lstlisting}[basicstyle=\ttfamily\footnotesize, frame=single]
pytest --nbval-lax --current-env notebooks
pytest --nbval-lax --current-env notebooks/07_deconvolution.ipynb
\end{lstlisting}
\texttt{--nbval-lax} lets a cell pass when it runs without an error; the
stored outputs are not compared, because plots, widget identifiers and
the last digits of noisy numbers change from run to run. A cell can still
be compared on purpose with the comment \texttt{\# NBVAL\_CHECK\_OUTPUT}.
\texttt{--current-env} runs the notebooks with the Python of the active
environment instead of a registered Jupyter kernel.

All 15 notebooks (00--09 and A--E, 132 code cells) pass in about two
minutes; most of the time is spent in the IQ simulation and correlation
of notebooks 03, 05 and~08. The notebook runs are therefore not part of
the plain \texttt{pytest} call (unit tests and doctests, a few seconds).
Files that the notebooks write (\texttt{notebooks/sim\_output/},
\texttt{notebooks/my\_example\_sky.json}) are ignored by git.

A certificate of the notebook runs is made in the same way as for the
unit tests (section~\ref{sec:certificate}):
\begin{lstlisting}[basicstyle=\ttfamily\footnotesize, frame=single]
pytest --nbval-lax --current-env notebooks --junitxml=logs/tests/notebooks.xml
python tools/testzertifikat.py --sprache en --titel "Notebook test certificate"
\end{lstlisting}

% ----------------------------------------------------------------
% 0.5.0: test certificate
\section{Test certificate}\label{sec:certificate}
\index{test certificate}

\texttt{tools/testzertifikat.py} turns a test log into a test
certificate. The tool uses only the Python standard library and imports
nothing from the package, so the same file serves other projects as well
(it comes from AAS7); version~1.1.0 added English output
(\texttt{--sprache en}; German stays the default) and the JUnit XML log
of pytest as default input.
\begin{lstlisting}[basicstyle=\ttfamily\footnotesize, frame=single]
pytest --junitxml=logs/tests/junit.xml
python tools/testzertifikat.py --sprache en
\end{lstlisting}
Without \texttt{--log} the tool takes the newest test log in
\texttt{logs/tests/} (an AAS7 JSON log \texttt{pytest\_*.json} first,
otherwise the newest \texttt{*.xml}). Further options:
\texttt{--min-tests N}, \texttt{--streng} (also no uncommitted state and no
unexpectedly passed test), \texttt{--titel}, \texttt{--aussteller},
\texttt{--ziel}, \texttt{--kein-pdf}, \texttt{--quiet}.

\paragraph*{Result.} PASSED when at least \texttt{--min-tests} tests
(default~1) ran and none failed or had an error in its setup; skipped
tests and expected failures do not count as failures. Exit code 0 PASSED,
1 FAILED, 2 no or unreadable test log.

\paragraph*{Storage and history.} Every certificate gets its own folder in
\texttt{certificates/}, which is never overwritten:
\texttt{TC\_<UTC time>Z\_<project>\_<branch>\_<commit>\_<checksum>}, the
last part being the first six characters of the SHA-256 checksum of the
test log. The folder contains \texttt{zertifikat.html} (readable without
further programs), \texttt{zertifikat.tex} and \texttt{zertifikat.pdf}
(if \texttt{pdflatex} is found), \texttt{zertifikat.json} (all data,
machine-readable) and an unchanged copy of the test log, whose checksum
is printed on the certificate. Each certificate adds a line to
\texttt{certificates/verlauf.csv}. The test logs in \texttt{logs/} are not
under version control, the certificates are.

\paragraph*{Tests.} \texttt{tests/test\_certificate.py} checks the tool
with small JUnit XML and AAS7 JSON logs: passed and failed certificates,
the checksum of the enclosed copy, \texttt{--min-tests} and
\texttt{--streng}, unique folder names and the history, the automatic
choice of the newest log, German as default and the exit codes of the
command line.

\section{Current tests}

% 0.3.0: table extended from two to all eight modules of tests/
% 0.3.1: longtable (may break across pages; as a tabular in a centre
%   environment the table was higher than the page)
\begin{longtable}{@{}l>{\raggedright\arraybackslash}p{8.4cm}@{}}
\toprule
File & Content \\
\midrule
\endfirsthead
\toprule
File & Content \\
\midrule
\endhead
\bottomrule
\endlastfoot
\texttt{tests/test\_import.py} & package version, import of all
  sub-packages, constants, all configuration files readable \\
\texttt{tests/test\_sky.py} & AAS6 file format and JSON round trip,
  demo field, flux conservation of the rendering, pixel positions,
  visibilities (analytic formulas, FFT of the rendered image, first null
  of the disc), map functions, command line tool \\
\texttt{tests/test\_instrument.py} & antenna configurations, horn and
  dish against notebook~B and \texttt{aas6sim}, primary beam at the
  HPBW, B210 profiles (inheritance, derived quantities, limits), SEFD and
  sensitivity formulas, ECEF and ENU coordinates, arrays of stations,
  instrument plots \\
\texttt{tests/test\_observation.py} & position of the Sun, transit
  time, hour angles and elevation limit, uv-track of an east-west
  baseline, visibility simulator against the sky model and against the IQ
  simulation, \texttt{VisData} round trip and selection, imaging
  (point source, phase centre shift, weightings, channel gridding),
  azimuth/elevation, parallactic angle and precession \\
\texttt{tests/test\_signal.py} & GMST, geometric delay, delay rate and
  $uvw$, fractional delay, quantisation (one bit, Van Vleck,
  efficiency), noise level, LO offset, station errors, SigMF output,
  fringe stopping, delay-rate search, FX and XF correlator, IQ imbalance \\
\texttt{tests/test\_solar.py} & flux of the Sun, solar transit at local
  noon, solar transit simulation (amplitude, solar rate), disc fit \\
\texttt{tests/test\_spectral.py} & velocity conversion, spectra of sky
  components and their JSON round trip, LSR correction, spectral
  visibility simulator (line at the LSR velocity, noise per channel),
  channel images \\
\texttt{tests/test\_deconvolution.py} & clean beam fit, H\"ogbom CLEAN
  (point sources, restored image, window, input checks), Wiener filter,
  reference image and RMS error, common deconvolver interface \\
\texttt{tests/test\_properties.py} & property-based tests with
  Hypothesis (section~\ref{sec:hypothesis}): spectra, sky components,
  geometry, receive path, sensitivity \\  % 0.4.0
docstrings in \texttt{src/interferosim/} & 13 doctests
  (section~\ref{sec:doctests}) \\  % 0.4.0
\texttt{tests/test\_certificate.py} & test certificate tool
  (section~\ref{sec:certificate}): passed/failed, checksum, criteria,
  unique folders and history, choice of the log, exit codes \\  % 0.5.0
\texttt{notebooks/*.ipynb} & all 15 notebooks run without errors (nbval,
  section~\ref{sec:nbval}; separate call) \\  % 0.6.0
\end{longtable}
